\chapterbackmatter{Solutions to Weinberg Exercises}
\ExerciseEditorialNote

\begin{enumerate}
\WeinbergSolution{1}
Write the center-of-mass energy as \(s=4\omega^2\), and define
\[
 \beta_e=\sqrt{1-\frac{m_e^2}{\omega^2}},
 \qquad
 \beta_\mu=\sqrt{1-\frac{m_\mu^2}{\omega^2}} .
\]
At tree level there is one annihilation graph.  With the overall phase
suppressed, its invariant amplitude is
\[
 \mathcal M=\frac{e^2}{s}\,
 \bar v(p_2)\gamma^\nu u(p_1)\,
 \bar u(k_1)\gamma_\nu v(k_2).
\]
The completeness relations reduce the unobserved-spin sum to two traces.
In the center-of-mass frame the result is
\[
 \overline{|\mathcal M|^2}
 =e^4\left[
  1+\frac{m_e^2+m_\mu^2}{\omega^2}
  +\beta_e^2\beta_\mu^2\cos^2\theta
 \right].
\]
The two-body phase-space formula therefore gives
\[
 \boxed{\;
 \frac{d\sigma}{d\Omega}
 =\frac{\alpha^2}{4s}\frac{\beta_\mu}{\beta_e}
 \left[
  1+\frac{4(m_e^2+m_\mu^2)}{s}
  +\left(1-\frac{4m_e^2}{s}\right)
   \left(1-\frac{4m_\mu^2}{s}\right)\cos^2\theta
 \right],\quad
 \alpha=\frac{e^2}{4\pi}.\;}
\]
This expression is zero below the muon threshold.  Angular integration
yields
\[
 \boxed{\;
 \sigma
 =\frac{4\pi\alpha^2}{3s}\frac{\beta_\mu}{\beta_e}
 \left[
  1+\frac{2(m_e^2+m_\mu^2)}{s}
  +\frac{4m_e^2m_\mu^2}{s^2}
 \right].\;}
\]
For \(m_e\to0\), this reduces to
\[
 \sigma=\frac{4\pi\alpha^2}{3s}\,
 \beta_\mu\left(1+\frac{2m_\mu^2}{s}\right).
\]

\WeinbergSolution{2}
Let
\[
 \pi\equiv\frac{\partial\mathcal L}{\partial\dot\Phi}
 =(D_0\Phi)^\dagger,\qquad
 \pi^\dagger\equiv\frac{\partial\mathcal L}{\partial\dot\Phi^\dagger}
 =D_0\Phi .
\]
The momentum conjugate to \(A^0\) vanishes.  Coulomb gauge and Gauss's law
form the electromagnetic constraints,
\[
 \boldsymbol\nabla\cdot\mathbf A=0,\qquad
 -\nabla^2A^0=\rho,\qquad
 \rho=iq\bigl(\Phi^\dagger\pi^\dagger-\pi\Phi\bigr).
\]
Thus \(A^0\) is auxiliary:
\[
 A^0(\mathbf x)=\int d^3y\,
 \frac{\rho(\mathbf y)}{4\pi|\mathbf x-\mathbf y|}.
\]
Writing \(\boldsymbol\Pi_\perp=\dot{\mathbf A}\), the nonzero independent
equal-time brackets are
\[
 [\Phi(\mathbf x),\pi(\mathbf y)]=i\delta^3(\mathbf x-\mathbf y),
 \qquad
 [A_i(\mathbf x),\Pi_{\perp j}(\mathbf y)]
 =i\delta^{T}_{ij}(\mathbf x-\mathbf y),
\]
where
\[
 \delta^T_{ij}(\mathbf x-\mathbf y)
 =\int\frac{d^3k}{(2\pi)^3}
 \left(\delta_{ij}-\frac{k_i k_j}{\mathbf k^2}\right)
 e^{i\mathbf k\cdot(\mathbf x-\mathbf y)} .
\]
The constrained Legendre transform gives
\[
 \begin{split}
 H={}&\int d^3x\left[
  \pi^\dagger\pi+(D_i\Phi)^\dagger(D_i\Phi)
  +m^2\Phi^\dagger\Phi+\lambda(\Phi^\dagger\Phi)^2
  +\frac12\boldsymbol\Pi_\perp^2+\frac12\mathbf B^2
 \right]\\
 &+\frac12\int d^3x\,d^3y\,
 \frac{\rho(\mathbf x)\rho(\mathbf y)}
 {4\pi|\mathbf x-\mathbf y|}.
 \end{split}
\]
This is expressed only in the requested canonical variables; \(A^0\) has
been eliminated.

For the interaction picture, let \(\phi,\phi^\dagger,\mathbf a\) be free
fields and set
\[
 \rho_I=iq\bigl(\phi^\dagger\dot\phi-\dot\phi^\dagger\phi\bigr).
\]
Expanding the spatial covariant derivatives then gives
\[
 \begin{split}
 H_I(t)={}&\int d^3x\left\{
 iq\,\mathbf a\cdot
 \bigl[\phi^\dagger\boldsymbol\nabla\phi
 -(\boldsymbol\nabla\phi^\dagger)\phi\bigr]
 +q^2\mathbf a^2\phi^\dagger\phi
 +\lambda(\phi^\dagger\phi)^2\right\}\\
 &+\frac12\int d^3x\,d^3y\,
 \frac{\rho_I(\mathbf x,t)\rho_I(\mathbf y,t)}
 {4\pi|\mathbf x-\mathbf y|}.
 \end{split}
\]
The first line supplies the scalar--scalar--photon vertex, the seagull
vertex, and the quartic matter vertex.  The second line is the instantaneous
Coulomb interaction.  As in Section 8.5, its noncovariant contribution
combines with transverse-photon exchange to reproduce the covariant photon
propagator.

\WeinbergSolution{3}
Let the initial scalar be at rest, let \(\omega\) and \(\omega'\) be the
initial and final photon energies, and let \(\theta\) be their angle.
Energy--momentum conservation gives
\[
 \omega'=\frac{\omega}
 {1+(\omega/m)(1-\cos\theta)}.
\]
There are two scalar-exchange graphs and the \(q^2\mathbf a^2\phi^\dagger
\phi\) seagull graph.  In a gauge with both photon polarizations orthogonal
to the initial scalar momentum, their gauge-invariant sum reduces to
\[
 \mathcal M=2q^2\,e^{\prime *}\cdot e .
\]
The result in another gauge is identical because replacing either
polarization by its photon momentum makes the full three-graph sum vanish.
For unobserved photon polarizations,
\[
 \frac12\sum_{e,e'}|e^{\prime *}\cdot e|^2
 =\frac12(1+\cos^2\theta).
\]
The laboratory two-body phase space therefore gives
\[
 \boxed{\;
 \frac{d\sigma}{d\Omega}
 =\frac{\alpha^2}{2m^2}
 \left(\frac{\omega'}{\omega}\right)^2
 (1+\cos^2\theta).\;}
\]
To integrate, put \(x=\omega/m\) and
\(\omega'/\omega=[1+x(1-\cos\theta)]^{-1}\).  Then
\[
 \int_{-1}^{1}dc\,
 \frac{1+c^2}{[1+x(1-c)]^2}
 =\frac{4(1+x)^2}{x^2(1+2x)}
 -\frac{2(1+x)}{x^3}\ln(1+2x).
\]
Hence
\[
 \boxed{\;
 \sigma=\frac{\pi\alpha^2}{m^2}
 \left[
 \frac{4(1+x)^2}{x^2(1+2x)}
 -\frac{2(1+x)}{x^3}\ln(1+2x)
 \right].\;}
\]
The bracket tends to \(8/3\) as \(x\to0\), so the low-energy limit is the
Thomson value \(8\pi\alpha^2/(3m^2)\).

\WeinbergSolution{4}
Let \(W_\mu\) be a complex vector field of charge \(q\), transforming as
\[
 W_\mu\longmapsto e^{iq\Lambda}W_\mu,\qquad
 A_\mu\longmapsto A_\mu+\partial_\mu\Lambda .
\]
Then
\[
 D_\mu W_\nu=(\partial_\mu-iqA_\mu)W_\nu,\qquad
 W_{\mu\nu}=D_\mu W_\nu-D_\nu W_\mu
\]
both transform covariantly.  One gauge-invariant answer is
\[
 \boxed{\quad
 \mathcal L=
 -\frac14F_{\mu\nu}F^{\mu\nu}
 -\frac12W_{\mu\nu}^\dagger W^{\mu\nu}
 -m^2W_\mu^\dagger W^\mu .
 \quad}
\]
Every Lorentz index is contracted between a field and its conjugate, so the
local phase cancels.  The mass term is therefore allowed even though a
photon mass would not be.

Gauge invariance alone does not make this answer unique.  For example the
Hermitian, dimension-four term
\[
 i\kappa qF_{\mu\nu}W^{\mu\dagger}W^\nu
\]
may be added with arbitrary real \(\kappa\), as may gauge-invariant
self-interactions.  The displayed minimal Lagrangian has the Proca
constraint and three physical polarizations.  A charged vector arising as
a gauge boson of a broken non-Abelian theory has the special magnetic
coupling conventionally corresponding to \(\kappa=1\).

\WeinbergSolution{5}
Label the process
\(e^-(p_1)e^-(p_2)\to e^-(p_3)e^-(p_4)\), and use
\[
 s=-(p_1+p_2)^2,\qquad
 t=-(p_1-p_3)^2,\qquad
 u=-(p_1-p_4)^2,\qquad s+t+u=4m^2 .
\]
The two final electrons can be attached in two ways.  Fermi statistics
fixes their relative minus sign:
\[
 \mathcal M=e^2\left[
 \frac{\bar u(p_3)\gamma^\mu u(p_1)\,
       \bar u(p_4)\gamma_\mu u(p_2)}{t}
 -\frac{\bar u(p_4)\gamma^\mu u(p_1)\,
       \bar u(p_3)\gamma_\mu u(p_2)}{u}
 \right].
\]
Spin completeness and the four-gamma trace give
\[
 \begin{split}
 \overline{|\mathcal M|^2}=2e^4\biggl[
 &\frac{s^2+u^2+8m^2(t-m^2)}{t^2}
 +\frac{s^2+t^2+8m^2(u-m^2)}{u^2}\\
 &+\frac{2(s^2-8m^2s+12m^4)}{tu}
 \biggr].
 \end{split}
\]
Because the final particles are identical, using the full solid angle
requires the factor \(1/2!\):
\[
 \boxed{\qquad
 \frac{d\sigma}{d\Omega}
 =\frac{\overline{|\mathcal M|^2}}{128\pi^2s}.
 \qquad}
\]
Equivalently, one may omit the factor and retain only one of the two
hemispheres.  As a useful check, for \(s\gg m^2\) the full-solid-angle
answer becomes
\[
 \boxed{\qquad
 \frac{d\sigma}{d\Omega}
 =\frac{\alpha^2}{2s}\,
 \frac{(3+\cos^2\theta)^2}{\sin^4\theta}.
 \qquad}
\]
It is symmetric under \(\theta\leftrightarrow\pi-\theta\), as required
when the two outgoing electrons cannot be distinguished.
\end{enumerate}
